\section{Connected clusters with repeated labels}
\label{sec:cluster-expansion}
Section~\ref{sec:ope-grouping} counted connected sets of distinct activity
labels. A logarithmic expansion is different: its vertices are numbered
occurrences, and the same activity label may occur several times. We now
prove the bound for this expansion, including its factorials and the
identity with the real logarithm of the gas. No positivity of the gas is
assumed in obtaining the identity.

Let $\Lambda$ be a finite periodic lattice. Each label $i$ has a nonempty
support $X_i$, and write $c(X)=C_B(X)$ and $h_B=(2B+1)^4$. We use
$|X|\le h_Bc(X)$ and covering subadditivity for intersecting supports.
The incompatibility relation $R$ is symmetric and reflexive, and
\[
                R(i,j)\quad\Longrightarrow\quad X_i\cap X_j\ne\varnothing.
\]
The converse is not required. Let $w_i\ge0$, $a\ge0$, $\gamma>0$,
$\delta>0$, and suppose
\begin{equation}\label{eq:cluster-vertex-input}
 \sup_{x\in\Lambda}\sum_{i:x\in X_i}w_i e^{(a+\delta)c(X_i)}\le\gamma,
 \qquad q_0=\frac{4h_B\gamma}{\delta}\le q<1.
\end{equation}
In particular $q>0$.

For an ordered tuple $z=(i_0,\ldots,i_m)$, let $T(z)$ be the graph
on $\{0,\ldots,m\}$ with an edge $\{r,s\}$ when $r\ne s$ and
$R(i_r,i_s)$. Define its connected coefficient and majorant by
\[
 u(z)=\sum_{\substack{G\subseteq T(z)\\G\ {\rm connected}}}(-1)^{|E(G)|},
 \qquad W(z)=\frac{|u(z)|}{(m+1)!}\prod_{r=0}^m w_{i_r}.
\]
For $m=0$ the graph with no edges is connected and $u(z)=1$.
Repeated labels are allowed in every tuple sum below.

\subsection{Cancellation before taking absolute values}
For any graph $T$ on $\{0,\ldots,m\}$,
\begin{equation}\label{eq:cluster-tree-graph}
 \left|\sum_{\substack{G\subseteq T\\G\ {\rm connected}}}
                      (-1)^{|E(G)|}\right|
                      \le\#\{\text{spanning trees of }T\}.
\end{equation}
Here is the cancellation argument. Fix a total order of the vertices
and root at zero. For each connected $G$, let $d(v)$ be its distance
from zero, and choose the least neighbor $p(v)$ at distance $d(v)-1$
as the parent of each $v\ne0$. The resulting parent edges form a
spanning tree $P$. They determine $d$ by their path lengths.

For a fixed tree $P$, define $U(P)\subseteq T$ to contain the tree
edges and all further edges satisfying these two rules: their levels
differ by at most one, and, if they join $v$ at level $r$ to $w$ at
level $r-1$, then $w\ge p(v)$ in the chosen order. Precisely the graphs
with
\[
                         P\subseteq G\subseteq U(P)
\]
have this breadth-first parent tree. One implication follows from
shortest-path distance and the least-parent rule. Conversely, the tree
paths show that a graph in this interval reaches every vertex in at
most $d(v)$ steps. The level restriction implies that no path from
zero reaches it in fewer steps, so the distances are exactly $d$.
The second rule makes $p(v)$ its least neighbor at the preceding
level. This proves the other implication and partitions all connected
graphs into these intervals.

The alternating sum over one interval is
$(-1)^m(1-1)^{|U(P)\setminus P|}$, hence has absolute value at most
one. Equivalently, if there is a free edge, toggling it pairs terms
with opposite signs. Summing over $P$ proves
\eqref{eq:cluster-tree-graph}. Taking absolute values separately for
all connected graphs would lose this cancellation and would not give
the required estimate.

\subsection{The weighted forest estimate}
A parent map $p:\{1,\ldots,m\}\to\{0,\ldots,m\}$ is called rooted
if iterating it always reaches the external root zero. Place a fixed
support $Y$ at zero and a support $X_{i_v}$ at vertex $v\ne0$.
Only assignments in which each child support meets its parent support
are retained. We claim
\begin{equation}\label{eq:cluster-forest-bound}
 \frac1{m!}\sum_{p\ {\rm rooted}}\sum_{i_1,\ldots,i_m}
 \mathbf1_{\{\text{all parent edges overlap}\}}
       \prod_{v=1}^m w_{i_v}e^{a c(X_{i_v})}
                            \le e^{\delta c(Y)}q^m.
\end{equation}
The labels in this sum need not be distinct.

Put $t=h_B\gamma/\delta>0$, $F(Y)=\max(\delta c(Y),1)$, and
introduce the nonnegative kernel
\[
 K(Y,i)=\frac{\mathbf1_{\{X_i\cap Y\ne\varnothing\}}
                        w_i e^{(a+\delta)c(X_i)}}{tF(Y)}.
\]
The site bound in \eqref{eq:cluster-vertex-input} implies
\[
 \sum_i\mathbf1_{\{X_i\cap Y\ne\varnothing\}}
                w_i e^{(a+\delta)c(X_i)}
       \le |Y|\gamma\le h_B\gamma c(Y)\le tF(Y).
\]
Thus every row of $K$ sums to at most one. For a fixed rooted parent
map, the sum over all label assignments of the product of its edge
kernels is at most one: sum the label at a leaf first, apply its row
bound, and repeat until no vertex remains. For infinite alphabets
one first restricts to a finite alphabet; the row bounds still hold,
and nonnegative convergence removes the restriction.

Let $d_v$ be the number of children of vertex $v$, including the
external root. The exponential series gives
\[
                 F(Z)^k\le k!e^{\delta c(Z)}\quad(k\ge0).
\]
Indeed if $\delta c(Z)\ge1$ this is one term in the exponential
series; otherwise use $F(Z)=1$ and $k!e^{\delta c(Z)}\ge1$.
Replacing each overlap weight with its kernel expression and collecting
the factors at parents gives
\begin{align*}
 &\prod_{v=1}^m\mathbf1_{\{X_{i_v}\text{ meets parent}\}}
                    w_{i_v}e^{a c(X_{i_v})}\\
 &\qquad\le t^m e^{\delta c(Y)}
                    \left(\prod_{v=0}^m d_v!\right)
                    \prod_{v=1}^m K(X_{i_{p(v)}},i_v),
\end{align*}
where $X_{i_0}=Y$. To check the reserve cancellation, before the
inequality there is a factor $e^{-\delta c(X_{i_v})}$ at every
nonroot vertex. Applying the bound on $F^{d_v}$ supplies exactly
one inverse factor there, leaving only $e^{\delta c(Y)}$ at the
external root.

It remains to sum the child factorials. Summing over all parent maps,
including those which are not rooted, can only increase the answer.
For fixed child counts $d_0+\cdots+d_m=m$ there are
$m!/\prod_vd_v!$ maps. Therefore
\[
 \sum_{p:\{1,\ldots,m\}\to\{0,\ldots,m\}}\prod_{v=0}^m d_v!
       =m!\binom{2m}{m}\le m!4^m.
\]
The binomial coefficient counts weak compositions of $m$ into $m+1$
parts. The inequality follows by comparing it with the sum of all
binomial coefficients of order $2m$. After division by $m!$, the
kernel bound and this count prove
\eqref{eq:cluster-forest-bound}, since $(4t)^m=q_0^m\le q^m$.
This also proves the case $m=0$, whose left side is one.

A distinguished root may carry a different nonnegative weight $v_i$.
If $\sum_i v_i e^{(a+\delta)c(X_i)}\le b$, multiply
\eqref{eq:cluster-forest-bound} with $Y=X_{i_0}$ by
$v_{i_0}e^{a c(X_{i_0})}$ and sum $i_0$. The result is the rooted
bound $bq^m$, with denominator $m!$ for the other vertices.

\subsection{A root anywhere in the tuple}
Let $Q$ be any property of an activity label and assume
\[
              \sum_{i:Q(i)}w_i e^{(a+\delta)c(X_i)}\le b.
\]
Apply \eqref{eq:cluster-tree-graph} to each occurrence graph and the
rooted forest bound with $v_i=\mathbf1_{Q(i)}w_i$ at a specified
vertex. This gives
\[
 \sum_{z}\frac{|u(z)|}{m!}\,
          \mathbf1_{Q(i_r)}\prod_{s=0}^m w_{i_s}e^{a c(X_{i_s})}
                           \le bq^m
\]
for every root position $r$. Relabeling occurrence positions preserves
both the connected coefficient and the tuple sum, so the same bound
holds for all $r$. If a tuple contains at least one label satisfying
$Q$, its cluster majorant, with denominator $(m+1)!$, is bounded by
$1/(m+1)$ times the sum of the displayed rooted integrands over all
$r$. Summing yields
\begin{equation}\label{eq:cluster-marked-size}
 \sum_{z:\,\exists r\,Q(i_r)}
 \frac{|u(z)|}{(m+1)!}\prod_{r=0}^m w_{i_r}e^{a c(X_{i_r})}
                           \le bq^m.
\end{equation}
This is the normalization that pays for choosing a root; no additional
factor $m+1$ remains.

If $u(z)\ne0$, its occurrence graph contains a connected spanning
subgraph, hence a spanning tree. Incompatibility edges imply support
overlap. Adding supports along this tree gives
\[
        c\left(\bigcup_{r=0}^m X_{i_r}\right)
                              \le\sum_{r=0}^m c(X_{i_r}).
\]
As $a\ge0$, the exponential weight of the actual union is therefore
bounded by the product weight in \eqref{eq:cluster-marked-size}.
Take $Q(i)$ to mean $x\in X_i$ and $b=\gamma$. The union contains
$x$ exactly when some label has this property. Finally sum over
$m\ge z$ to obtain
\begin{equation}\label{eq:cluster-local-tail}
 \sum_{m\ge z}\sum_{i_0,\ldots,i_m}
 \mathbf1_{\{x\in\bigcup_rX_{i_r}\}}
      W(i_0,\ldots,i_m)e^{a c(\bigcup_rX_{i_r})}
                       \le\frac{\gamma q^z}{1-q}.
\end{equation}
Thus the connected estimate \eqref{eq:leaf-cluster-tail} follows
from a vertex site bound and the explicit smallness inequality in
\eqref{eq:cluster-vertex-input}: use the exponent $a+e$ there in
place of $a$, the cluster vertex price $t$ in place of $\gamma$,
and the tail parameter $\gamma$ of that section in place of $q$.
The reserve $\delta$ pays for tree branching and is distinct from the
reserve $e$ used later for spatial derivative weights.

\subsection{The exact logarithm identity}
Now let $v_i\in\mathbb R$ be possibly signed activities and apply the
preceding estimates to $w_i=|v_i|$. Define the gas and connected sum
\[
 Z=\sum_{A\ {\rm finite,compatible}}\prod_{i\in A}v_i,
 \qquad
 \mathcal L=\sum_{m\ge0}\frac1{(m+1)!}
         \sum_{i_0,\ldots,i_m}u(i_0,\ldots,i_m)\prod_{r=0}^m v_{i_r}.
\]
A finite set is compatible if it contains no incompatible distinct
labels. The empty set contributes one. We will prove
\begin{equation}\label{eq:cluster-exact-log}
                         Z=e^{\mathcal L}>0,\qquad \log Z=\mathcal L.
\end{equation}
First both series converge absolutely. Nonempty supports and
\eqref{eq:cluster-vertex-input} imply $\sum_i|v_i|\le|\Lambda|\gamma$,
so the gas is dominated by $\prod_i(1+|v_i|)\le e^{|\Lambda|\gamma}$.
For the logarithm, use \eqref{eq:cluster-marked-size} with $Q$ always
true and $b=|\Lambda|\gamma$, and drop the support weights. Its
absolute sum at order $m$ is at most $|\Lambda|\gamma q^m$; hence
the total is at most $|\Lambda|\gamma/(1-q)$.

Begin with a finite alphabet $D$. Let $P(t)$ be its gas polynomial,
with each activity replaced by $tv_i$, and let
\[
                 L(t)=\sum_{m\ge0}\ell_m t^{m+1},\qquad
 \ell_m=\frac1{(m+1)!}\sum_{i_0,\ldots,i_m\in D}
                                    u(i_0,\ldots,i_m)\prod_rv_{i_r}.
\]
The bound $|\ell_m|\le |\Lambda|\gamma q^m$ shows that $L$ is
differentiable for $|t|<q^{-1}$ by termwise differentiation. Indeed
on $|t|\le u<q^{-1}$ the derivative terms are bounded by the summable
series $|\Lambda|\gamma(m+1)(qu)^m$. This interval contains $[0,1]$.

We verify the coefficient identity $P'=L'P$, rather than assuming a
formal logarithm can be evaluated. On a tuple of $n$ occurrences,
the sum of $(-1)^{|E(G)|}$ over all subgraphs of its incompatibility
graph is $(1-1)^{|E(T)|}$. It is one if the graph has no edge and zero
otherwise. Because $R$ is reflexive, a tuple with repeated labels
has an edge between equal occurrences and contributes zero. Thus
summing this all-graph weight over tuples, multiplying by their
activities, and dividing by $n!$ gives exactly the coefficient $p_n$
of $P(t)$: each compatible $n$-element set has $n!$ enumerations.

For $n=m+1$, distinguish occurrence zero in an arbitrary graph. Its
connected component consists of zero and some $k$ of the other $m$
vertices. Choose those vertices in $\binom{m}{k}$ ways. A graph is
then uniquely specified by a connected graph on this component and
an arbitrary graph on its complement, with no cross edges. Edge
signs and label products factor. Summing the choices and dividing
by $m!$ gives the exact recurrence
\[
                  (m+1)p_{m+1}
                       =\sum_{k=0}^m(k+1)\ell_k p_{m-k}.
\]
Absolute convergence of $L'$ and the finiteness of $P$ allow its
coefficients to be multiplied and summed, proving $P'=L'P$ on the
interval. Hence
\[
                 \frac{d}{dt}\{P(t)e^{-L(t)}\}=0.
\]
At zero, $P(0)=1$ and $L(0)=0$. The mean value theorem implies
$P(t)e^{-L(t)}=1$ throughout the interval, and evaluation at $t=1$
proves \eqref{eq:cluster-exact-log} for the finite alphabet.

Finally let finite alphabets increase to all labels. Every gas term
and every occurrence tuple is eventually retained. The absolute
majorants already proved are independent of this restriction.
Dominated convergence therefore gives $P_D(1)\to Z$ and
$L_D(1)\to\mathcal L$. Continuity of the exponential yields
$Z=e^{\mathcal L}$, and its strict positivity gives the claimed
real logarithm. For an arbitrary index set the same argument uses
the directed family of all finite subsets; alternatively absolute
summability leaves only countably many nonzero terms. This completes
the proof, with no assumption of positivity or a choice of a complex
logarithm branch.

Applied to the uniformly dominated tilted activities of
Section~\ref{sec:leaf-expansion}, this result constructs their positive
normalization and logarithmic cluster series once the vertex site
bound \eqref{eq:cluster-vertex-input} is established. The joint bounds
proved there then justify differentiating it. Constructing the actual
block-integral activities and proving their vertex bounds remain
separate analytic tasks.
